\documentclass[11pt]{article}
\usepackage[utf8]{inputenc}
\usepackage[T1]{fontenc}
\usepackage{lmodern}
\usepackage[margin=1in]{geometry}
\usepackage{csquotes}
\usepackage[style=authoryear,backend=biber,sorting=nyt,maxbibnames=99,giveninits=true,uniquename=init]{biblatex}
\usepackage[hidelinks]{hyperref}
\addbibresource{references.bib}
\begin{document}
\section*{Longitudinal Data Gaps and Monitoring Imperatives}

\subsection{Overview}

A recurring message across the subsurface (tile) drainage literature is that the current evidence base remains provisional. Many headline results derive from short records, single sites, or narrow agro-ecological conditions, and their authors explicitly caution against premature generalisation. The dominant frontier signal is therefore not a settled conclusion but a research agenda: the field is asked to extend monitoring in time, replicate it across space, validate models and emerging technologies at field scale, harmonise methods so that accumulating data can be synthesised, deepen mechanistic understanding of transport pathways and process interactions, and connect environmental endpoints to agronomic, economic, and social outcomes and to decision-making. Several authors pitch these needs at the level of the whole discipline. \textcite{Rogger2017}, for example, argue for complex-systems thinking that links processes across time scales, long-term experiments on physical--chemical--biological process interactions, a focus on connectivity and patterns across spatial scales, typologies and meta-analyses that generalise findings, multidisciplinary consortia with shared data and terminology, and long-term observation under controlled or known boundary conditions wherever possible. This section organises the specific monitoring and research needs identified across the reviewed studies along temporal, spatial, methodological, mechanistic, integrative, and programmatic dimensions.

The scale of the outstanding uncertainty is captured by the sheer frequency of direct appeals for more work. Studies conclude, variously, that ``more research is needed'' and that ``more verification would be beneficial'' \parencite{SHALIT1995}, that ``the complexity of this subject requires more research'' \parencite{Malone2014}, that ``further research in this area is needed'' \parencite{Woo2023}, and that ``further evaluations are warranted before any concrete conclusions could be drawn'' \parencite{GUPTA1993}. These are less rhetorical closings than signposts to genuine gaps in evidence, method, and synthesis, which the subsections below make explicit.

\subsection{Extending the Temporal Record}

\subsubsection{The Case for Long-Term, Multi-Year Monitoring}

The single most widely voiced need is for records long enough to separate genuine treatment effects from the natural variability of weather, hydrology, and crop growth. \textcite{Clifford2023} emphasise the need for truly long-term experiments and monitoring to assess management impacts accurately, and \textcite{BUHLER1993} caution that a lag of several years may be necessary before the effect of an altered practice can be objectively evaluated. \textcite{Ali2013} judge that longer monitoring periods---at least five years---and continued data collection are desirable, warning that long-term consequences may require continued observation if drainage systems are expanded in length or depth over time. In Egypt, \textcite{Ali2001} report that monitoring will continue and that longer time series for the whole country, expected within a few years, should enable more detailed analysis as the records lengthen.

At the site scale, many studies attach their conclusions to explicit requirements for additional years of observation. \textcite{Jaynes2008} note that continued monitoring of a denitrification wall over additional years is necessary to quantify its effectiveness against natural variability, to establish long-term cost per unit of nitrate removed, and to gauge its potential as a management tool. \textcite{Dalhem2019} expect continued sampling to reveal how long the positive effects of their treatments persist, and \textcite{Osterholm2005} judge it worthwhile to monitor a test site for several more years to assess the functionality of a combined technique. Similar calls for verification over longer periods come from \textcite{ROLLIN1987}, who plan observations spanning a longer service life; \textcite{Wesström2015}; \textcite{Reinhart2020}, who seek more widespread and longer-duration evaluations of a practice's potential impact; \textcite{Rijal2012}, who stress the need to assess the replicability of results in time; and \textcite{GiménezValero2025}, who conclude that additional multi-year measurements are needed to confirm whether observed differences persist under comparable conditions. Underlying all of these is the disciplinary argument of \textcite{SPALING1995} for long-term monitoring of environmental change at various scales.

\subsubsection{Temporal Resolution and Seasonal Coverage}

Record length is necessary but not sufficient: several authors highlight gaps in temporal \emph{resolution} and seasonal coverage. \textcite{Klaiber2020} call for additional years of monitoring together with improved, year-round measurement---especially during the non-growing season---to understand seasonal phosphorus exports and reach more robust conclusions. \textcite{Hodaj2017} stress avoiding gaps in high-frequency records, advising that future studies check instruments as frequently as possible. \textcite{Wesström2015} add a timing consideration specific to infrastructure interventions, noting that it is more appropriate to monitor nutrient leaching after drain repair.

\subsection{Broadening Spatial and Geographic Scope}

A second major frontier is testing whether site-specific findings hold across other locations, soils, climates, and management regimes. \textcite{Menegaz2025} caution that additional studies are needed before broad adoption of their results can be encouraged beyond local conditions, and \textcite{Zimmerman2017b} call for similar data from other sites and years. \textcite{Ali2001} plan to collect data from further sample areas reflecting Egypt's diverse agro-ecological conditions. \textcite{Hodaj2017} argue that exploring the two-stage ditch in other settings would benefit the practice, and \textcite{Bigelow2001} recommend research into effects where putting greens are severely compacted, have low percolation rates, or are maintained in unfavourable environments such as shade or excessive wetness, while suggesting the technology may have merit under other situations. \textcite{Terrado2015} likewise conclude that further monitoring and research in support of their findings is warranted. At the level of evidence synthesis, \textcite{Wu2025} argue that future syntheses require more robust, geographically diverse, and methodologically harmonised studies, and that a broader experimental base will let future meta-analyses better disentangle management effects.

\subsection{Field-Scale Validation of Models and Emerging Technologies}

Much reviewed work rests on models, pilot studies, or laboratory-scale results whose authors defer confident application until field-scale validation is available. \textcite{Mitchell2025} state plainly that ``field-scale research is needed to validate these results'' and that ``further research to validate our models is needed before they are widely applied.'' \textcite{Kalcic2018} similarly conclude that continued efforts to validate practice effectiveness at field and reach scales will be needed to verify whether models are correct, marking these as areas ripe for future development. \textcite{Garmdareh2018} recommend more rigorous testing of a model before its extensive use, and \textcite{Yu2016} call for validating their conclusions with more experiments. \textcite{SHALIT1995} identify a specific target, arguing that further investigation should determine how well the leaching-fraction model predicts concentrations arising from preferential flow. For emerging technologies, \textcite{DarziNaftchali2018b} stress that environmental effects must be quantified at the pilot-area scale before large-scale implementation, and \textcite{GUPTA1993} require further measurements and evaluations before concrete conclusions can be drawn.

\subsection{Monitoring Infrastructure, Instrumentation, and Data Tools}

A distinct set of needs concerns the practical apparatus of monitoring---the samplers, loggers, diagnostic methods, and software that make long-term data collection feasible and reliable. \textcite{Kanwar1999} describe how automatic samplers can collect discrete water samples alongside continuous ones, how additional equipment can be connected to data loggers to measure soil temperature or moisture, and how relays or pressure switches could replace electronic transmitters to detect sump-pump operation; they advise minimising the distance between equipment and loggers, downloading data automatically to reduce loss, and performing annual maintenance including cleaning check valves and calibrating flowmeters. Diagnostic methods are themselves a research target: \textcite{Petersen2012} propose reducing pressure variations by doubling the equipment and injecting smoke from both ends of a drain section---potentially enabling work with longer drain lines---while noting that further study is needed to resolve transport pathways and the reasons for observed clustering. \textcite{MartínezBeltrán2007} recommend observing selected pipes to determine the state of the envelope and perforations, observing the surrounding soil for its influence on radial flow in the trench, and permanently monitoring the sampling process to verify that favourable results are maintained over time.

Software and data infrastructure form an increasingly prominent part of this frontier. \textcite{Shedekar2018} advise readers to explore data products and tools further, to consult the National Elevation Dataset and USGS resources for the latest availability, and to track ongoing developments, while arguing that improved drainage-management practices are needed to increase adaptation to a changing climate. A series of tool-development efforts describe themselves as explicitly ongoing and open to community input: \textcite{Falasy2025} note continuing development---spanning topologically sound, one-way-flow tile networks and application tools for configuring the burying and sizing of drainage networks---and welcome testing with additional datasets; \textcite{Falasy2025b} report ongoing revision based on test-user feedback and encourage application beyond the samples provided; and \textcite{Falasy2024} state that development of their tools will be continual and invite testing with other data.

\subsection{Standardisation, Harmonisation, and Comparability}

For the accumulating body of monitoring data to support robust synthesis, several authors identify a need for standardised, harmonised methods. \textcite{Wu2025} tie the value of future meta-analyses to methodologically harmonised studies and to comprehensive, long-term monitoring of the full suite of environmental impacts. \textcite{Zhang2018} set out a detailed programme for environmental indicators: improving the quality and comparability of existing indicators through harmonised data collection, monitoring, interpretation, and reporting; developing new indicators for emerging environmental and policy concerns while ensuring scientific credibility, sensitivity to change, and ease of interpretation; establishing standardised evaluation procedures, including standards, thresholds, reference points, and complementary validation information; recognising data gaps and improving methods for upscaling, uncertainty assessment, and analysis of contradictory indicators; and adapting environmental models for indicator modelling while incorporating earth-observation data and reducing barriers to remote-sensing access. These priorities complement the call by \textcite{Rogger2017} for typologies, meta-analyses, and shared data and terminology across research consortia.

\subsection{Deepening Mechanistic Understanding}

A further frontier concerns mechanistic understanding: identifying the pathways, relationships, and process interactions that current data only partially resolve. \textcite{Malone2014} conclude that the complexity of the subject requires more research, and \textcite{Cui2016} state that more research is needed to define the relationships they observed. \textcite{Petersen2012} identify unresolved transport pathways and clustering patterns requiring further investigation. \textcite{Tomer2008} intend to pursue their central idea as a research avenue, to undertake bacterial source tracking within their watershed, and note that a specific transport pathway remains a subject for future research. \textcite{Ali2013} call for a separate geochemical modelling study to test whether geochemical reactions bind free heavy metals to the solid matrix. Process-level questions extend to plant response: \textcite{GiménezValero2025} note that an observation consistent with the role of phosphorus in flowering merits further study under controlled conditions, and \textcite{Tamura2021} report that further experiments are required to clarify a specific point. \textcite{LAL1993} indicate that plant-nutrient results will appear in a companion report and that a long-term experiment is being analysed to test their hypotheses. These needs echo the overarching argument of \textcite{Rogger2017} for long-term experiments on physical--chemical--biological process interactions and for attention to connectivity and patterns across spatial scales.

\subsection{Integrating Environmental, Agronomic, Economic, and Social Evidence}

Several authors argue that monitoring must move beyond environmental endpoints to encompass agronomic, economic, and social dimensions and to connect with decision-making. \textcite{Zhang2018} call for improved linkage of environmental indicators with those tracking economic and social change, in terms of both causes and effects. \textcite{Zhu2025} require further assessment of actual growth in conjunction with economic, social, and environmental factors, and \textcite{Brikmans2025} argue that future work should elucidate strategies for prolonging physical benefits while critically assessing the agronomic and economic feasibility of such approaches. \textcite{Monaghan2016} identify a need for research on the practical implications of management changes---in their case, keeping cows off pasture for extended periods. At the science--policy interface, \textcite{Letey1999} argue that plans should be updated as monitoring, additional studies, and local actions reveal new facts, and that the scientific community should evaluate its role in policymaking and focus research on questions of greatest value to decision makers as well as to scientific peers. \textcite{Clifford2023} call for scientific assessment and prioritisation of which approaches have the largest effects under what circumstances, alongside truly long-term experiments. \textcite{Kalcic2018} likewise argue that further work is needed to assess current knowledge and prioritise knowledge gaps, and that clear messaging through education and extension will be imperative to achieve implementation at the scales needed to address water-quality needs in the Mississippi River Basin and Great Lakes regions.

\subsection{Ongoing and Planned Monitoring Programmes}

Finally, the literature documents monitoring programmes that are explicitly ongoing or planned, underscoring that many conclusions are interim by design. \textcite{Tournebize2012} report that a performance-assessment programme was established in 2011 to determine whether co-constructed wetlands comply with new regulations, that intensive hydrological monitoring will continue, and that resizing will be considered if needed. \textcite{Kjaersgaard2016} note that monitoring was to begin in spring 2016, and \textcite{Ali2001} confirm that monitoring will continue. \textcite{Ohlendorf2002} describe how monitoring data were used to identify needs for contingency plans and to recommend further research and monitoring to improve site management, with additional studies subsequently conducted to resolve unknowns flagged by a risk assessment. Other authors signal forthcoming results: \textcite{Drury2014} indicate that data from remaining treatments will appear in a subsequent paper. Further calls for continued or expanded empirical effort come from \textcite{Rekha2011}, who state that more treatment studies are required, and from \textcite{Ogino2007}, who hope that research and development in the field will be actively promoted. Collectively, these programmes and appeals reinforce the central frontier message: robust, generalisable conclusions about subsurface drainage await the longer, broader, better-harmonised, and more integrative evidence base that the field is actively working to build.

\printbibliography
\end{document}
